\section{Performance Evaluation}
\textbf{What we measure.} Each step is measured with the score that fits its task. \emph{Finding which clauses are present} is a yes/no decision for each of the 4{,}182 (contract, category) test pairs. We score it with micro-F1 (every decision counts the same), macro-F1 (every category counts the same), and the accuracy, precision, recall and specificity from the confusion matrix. \emph{Finding the clause text} is scored with token-level F1, exact match, and the share of answers that overlap the real clause with F1~$\geq$~0.5. \emph{Summarizing} is scored with ROUGE-L against two different sets of reference sentences. The \emph{whole system} is scored by the share of real clauses that get through every step, and the risk layer by how many clauses it finds at each risk level. Speed is measured as the time and memory the app needs (Table~\ref{tab:appbench}).

\textbf{How we test.} Every model is tested once on the same 102 test contracts, using a separate notebook that loads the saved models from disk. Uncertainty is measured with a bootstrap over contracts, and we check whether the main result changes with the amount of training data, the random seed, the way window scores are combined, the decision threshold and the train/test split.

Everything in this chapter is measured on the 20\% test set: the same 102 unseen contracts for every model. All numbers come from the October 2026 re-training (Chapter~4) unless a run is marked as from August 2026. They were produced by a separate test notebook that loads the trained models from disk, so nothing left over from training can affect them.
\subsection{Presence Classification}
\subsubsection{Model comparison}
Table~\ref{tab:ensemble} and Figure~\ref{fig:ensemble} compare all the presence models on the 4{,}182 (contract, category) test decisions. Two things stand out. First, the window-based transformer \emph{loses on its own} to the simple TF--IDF model that reads the whole document (0.692 vs.\ 0.775). Legal wording is distinctive enough that word evidence from the whole document beats meaning-based evidence from small windows, and taking the maximum window score creates too many false alarms, just as Section~\ref{sec:mil} predicted. Second, the AND-ensemble, which reports a clause only when \emph{both} models agree, removes many of each model's false alarms and comes out on top (0.779). The lead over the simple model is small on 102 test contracts, so we do not rely on it: Section~\ref{sec:bootstrap} measures exactly how small it is and finds that this gap cannot be told apart from zero. The conclusion we stand behind is about method: \emph{simple word-based models are strong at deciding which clauses a contract contains, and the transformer is useful as a filter that removes false alarms, not as a replacement}.

\begin{table}[!htbp]
\centering
\caption{Presence model comparison}
\label{tab:ensemble}
\begin{adjustbox}{max width=\textwidth}
\begin{tabular}{lc}
\toprule
\textbf{Model} & \textbf{micro-F1}\\
\midrule
\textbf{Ensemble AND (both must agree)} & \textbf{0.779}\\
TF--IDF baseline (full document)        & 0.775\\
Ensemble AVG (mean probability)         & 0.718\\
Transformer (max-pool) alone            & 0.692\\
Ensemble OR (either fires)              & 0.692\\
\addlinespace
\multicolumn{2}{l}{\emph{August 2026 control run, Section~\ref{sec:budget}:}}\\
Transformer trained on all 53{,}662 windows & 0.694\\
Ensemble AND on that transformer            & 0.782\\
\bottomrule
\end{tabular}
\end{adjustbox}

{\footnotesize\singlespacing \emph{Note.} micro-F1 over all 4{,}182 test decisions. The transformer (0.6924) is just ahead of OR (0.6916).\par}
\end{table}

\begin{figure}[!htbp]
\centering
\includegraphics[width=0.75\textwidth]{ensemble_comparison.png}
\caption{Presence model micro-F1}
\label{fig:ensemble}
\end{figure}
\subsubsection{Confusion-matrix analysis}
Our test plan required a full confusion-matrix analysis of the final model. Table~\ref{tab:cm} gives the matrix for the AND-ensemble over all test decisions, Figure~\ref{fig:cm} shows the same matrix as a heatmap, and Table~\ref{tab:metrics} lists every score taken from it.

\begin{table}[!htbp]
\centering
\caption{AND-ensemble confusion matrix}
\label{tab:cm}
\begin{adjustbox}{max width=\textwidth}
\begin{tabular}{lcc}
\toprule
 & \textbf{Predicted Absent} & \textbf{Predicted Present}\\
\midrule
\textbf{Actually Absent}  & TN = 2{,}525 & FP = 309\\
\textbf{Actually Present} & FN = 291     & TP = 1{,}057\\
\bottomrule
\end{tabular}
\end{adjustbox}
\end{table}

\begin{table}[!htbp]
\centering
\caption{Confusion-matrix metrics}
\label{tab:metrics}
\begin{adjustbox}{max width=\textwidth}
\begin{tabular}{llc}
\toprule
\textbf{Metric} & \textbf{Formula} & \textbf{Value}\\
\midrule
Accuracy    & $(TP+TN)/\text{all}$          & 85.65\%\\
Precision   & $TP/(TP+FP)$                  & 77.38\%\\
Recall (Sensitivity) & $TP/(TP+FN)$         & 78.41\%\\
Specificity & $TN/(TN+FP)$                  & 89.10\%\\
F1 score    & $2PR/(P+R)$                   & 0.7789\\
\bottomrule
\end{tabular}
\end{adjustbox}
\end{table}

\begin{figure}[!htbp]
\centering
\includegraphics[width=0.6\textwidth]{confusion_matrix.png}
\caption{Confusion matrix heatmap}
\label{fig:cm}
\end{figure}

Precision and recall are balanced (77.4\% and 78.4\%), which shows the AND rule did its job: taking the maximum window score alone loses precision, and requiring both models to agree wins it back without losing recall. The specificity of 89.1\% means the system rarely claims a clause that is not there; for a legal-review tool, that is the more harmful kind of mistake.
\subsubsection{Per-category behaviour}
Figure~\ref{fig:percat} shows the per-category F1 for the strongest categories. The common, standard categories are almost solved (\emph{Document Name} 1.000, \emph{Parties} 0.995, \emph{Agreement Date} 0.963, \emph{Governing Law} 0.931). At the other end, five categories with at most seven test examples score an F1 of 0.000 with the AND-ensemble. Table~\ref{tab:zerof1} shows exactly what this means. In every case the ensemble marked the category as present in \emph{no} test contract (no true and no false positives), so F1 = 0.000 is correct: none of the 2--7 real clauses was found. But precision is undefined ($0/0$) rather than zero. The transformer on its own does find some of these clauses; the AND rule removes them because the TF--IDF model never agrees (Section~\ref{sec:raretail}). With only two to seven examples per category, these scores are not reliable either way. This weakness comes from the dataset, not from any particular model, and it is why our main results use micro-averaged rather than macro-averaged scores.

\begin{table}[!htbp]
\centering
\caption{Categories with zero F1}
\label{tab:zerof1}
\small
\begin{adjustbox}{max width=\textwidth}
\setlength{\tabcolsep}{4pt}
\begin{tabular}{>{\raggedright\arraybackslash}p{3.7cm}cccc c}
\toprule
\textbf{Category} & \shortstack{\textbf{Test}\\\textbf{examples}} & \shortstack{\textbf{AND}\\\textbf{TP / FN}} & \shortstack{\textbf{AND}\\\textbf{precision}} & \shortstack{\textbf{Transformer alone}\\\textbf{F1 (TP/FP/FN)}} & \shortstack{\textbf{TF--IDF}\\\textbf{F1}}\\
\midrule
Most Favored Nation              & 7 & 0 / 7 & n/a & 0.000 (0/0/7) & 0.000\\
No-Solicit of Customers          & 6 & 0 / 6 & n/a & 0.323 (5/20/1) & 0.000\\
Price Restrictions               & 4 & 0 / 4 & n/a & 0.087 (1/18/3) & 0.000\\
Source Code Escrow               & 2 & 0 / 2 & n/a & 0.000 (0/0/2) & 0.000\\
Unlimited\slash All-You-Can-Eat License & 2 & 0 / 2 & n/a & 0.500 (1/1/1) & 0.000\\
\bottomrule
\end{tabular}
\end{adjustbox}

{\footnotesize\singlespacing \emph{Note.} n/a: the model never said ``present'', so precision is $0/0$ and has no value. Threshold 0.5, October 2026 models.\par}
\end{table}

\begin{figure}[!htbp]
\centering
\includegraphics[width=0.85\textwidth]{per_category_f1.png}
\caption{Top-12 per-category F1}
\label{fig:percat}
\end{figure}
\subsection{Span Extraction}
\label{sec:span}
We tested the question-answering model on all 2{,}667 test windows that contain an answer (Table~\ref{tab:span}). A token-F1 of 0.764, with 82.8\% of answers overlapping the real clause at F1 $\geq$ 0.5, means that when the clause is in front of the model, it captures most of the right text about eight times out of ten. The lower exact-match score did not surprise us: marked legal clauses are long and their exact start and end are often debatable (should the section number at the start be included?), so overlap is the more useful score here. For comparison, the best published CUAD result uses DeBERTa-xlarge, which has about ten times more parameters than our 67M-parameter model.

Section~\ref{sec:bootstrap} argues that a score on 102 contracts should always be read together with its confidence interval, and that is just as true here. So we ran the \emph{same} bootstrap over contracts. Here the clauses are even more linked to each other than in the presence task: the 2{,}667 span decisions come from only 102 contracts, and a contract with unusual writing (numbered sub-clauses, many defined terms) moves all of its clauses in the same direction. So Table~\ref{tab:span} gives two intervals for each score: one that resamples \emph{contracts} (the correct way) and one that resamples \emph{clauses} one by one (the naive way), because the difference between them is itself useful.

\begin{table}[!htbp]
\centering
\caption{Span extraction results}
\label{tab:span}
\small
\begin{adjustbox}{max width=\textwidth}
\begin{tabular}{lccc}
\toprule
\textbf{Metric} & \textbf{Score} & \textbf{95\% CI (by contract)} & \textbf{95\% CI (by clause)}\\
\midrule
Token-level F1                & 0.764 & $[0.743,\ 0.782]$ & $[0.752,\ 0.775]$\\
Overlap (token-F1 $\geq$ 0.5) & 82.8\% & $[80.0\%,\ 85.2\%]$ & $[81.4\%,\ 84.2\%]$\\
Exact Match                   & 35.5\% & $[32.6\%,\ 38.6\%]$ & $[33.7\%,\ 37.3\%]$\\
\bottomrule
\end{tabular}
\end{adjustbox}

{\footnotesize\singlespacing \emph{Note.} 95\% bootstrap intervals, 10{,}000 resamples. Quote the interval by contract. The interval by clause is shown only to make clear how much it understates the uncertainty.\par}
\end{table}

The interval for token-F1 when resampling contracts is $[0.743, 0.782]$, roughly \emph{twice} as wide as the $[0.752, 0.775]$ we would get by treating clauses as independent. This difference is the real cost of ignoring that clauses from the same contract are linked, and it is the same effect that made the contract-level bootstrap necessary in Section~\ref{sec:bootstrap}. It also shows how precise this number is: differences smaller than about two points of token-F1 cannot be measured on a test set of this size, which is worth remembering before comparing 0.764 with any published figure.

\textbf{Limit of this test.} These numbers measure extraction \emph{when the window contains the answer}. More exactly, each test window is the 1{,}200-character focus window placed around the real answer (Section~\ref{sec:leakage}), so the model is indirectly told where the answer is. This makes it a best-case score, not a score the real system can reach. The full-document test in the CUAD paper (precision at 80\% recall over all windows, including those with no answer) needs a heavier ranking test that we did not run, so we make no claim about that score. Section~\ref{sec:e2e} measures the other side: what happens to extraction when the window is chosen by the presence step instead of being given.
\subsection{Summarization}
\label{sec:summ}

\subsubsection{What the reference set actually is}
Before reading any number in this section, one fact about the test must be stated clearly, because it decides what the score can mean. CUAD has no plain-English versions of clauses, so there was no ready-made set of target sentences to train or test with. So we \textbf{wrote the targets ourselves}: 41 one-sentence summaries, one per category, written from the weaker party's side (Section~\ref{sec:summmethod}). Every training pair and every test pair uses one of these 41 sentences. The targets do not change from clause to clause at all: two \emph{Cap on Liability} clauses with completely different limits have exactly the same target sentence.

This design choice has an effect we should admit openly: with 41 different targets for 2{,}667 test clauses, a model can score well just by \emph{recognising} the category and repeating its sentence, and the score cannot tell this apart from real summarizing. The test below is designed to separate the two.

\subsubsection{Two reference sets, the same 150 clauses}
We test the fine-tuned FLAN-T5-small on 150 test clauses (chosen with seed 42) against two different sets of reference sentences:

\begin{enumerate}[itemsep=2pt]
  \item \textbf{Template references}: the category template, exactly as in training. This is how the main score was produced.
  \item \textbf{Clause-specific references}: 150 one-sentence summaries, one per clause, written to describe \emph{that} clause, with its actual limit, its actual notice period and its actual parties. They were drafted with the help of a large language model working only from the clause text, without seeing any model output, and they are fully separate from the 41 templates. They are a trial reference set, not written by a lawyer, and we claim nothing more for them.
\end{enumerate}

Because the clauses are the same in both tests, the difference between the two scores shows exactly one thing: how much of the ROUGE-L score comes from the reference sentences all being the same, rather than from real summarizing. Table~\ref{tab:sumeval} shows the result. The October 2026 re-trained summarizer gives exactly the same scores as the original.

\begin{table}[!htbp]
\centering
\caption{Summarizer ROUGE-L}
\label{tab:sumeval}
\small
\begin{adjustbox}{max width=\textwidth}
\begin{tabular}{llcc}
\toprule
\textbf{System} & \textbf{Reference set} & \textbf{ROUGE-L} & \textbf{95\% CI}\\
\midrule
Fine-tuned FLAN-T5-small & 41 category templates    & 0.775 & $[0.718,\ 0.833]$\\
Fine-tuned FLAN-T5-small & clause-specific (pilot)  & 0.116 & $[0.105,\ 0.128]$\\
Zero-shot FLAN-T5-small  & clause-specific (pilot)  & 0.299 & $[0.275,\ 0.322]$\\
\midrule
\emph{Category template itself} & clause-specific (pilot) & 0.115 & $[0.104,\ 0.127]$\\
\bottomrule
\end{tabular}
\end{adjustbox}

{\footnotesize\singlespacing \emph{Note.} 95\% bootstrap intervals over the 150 clauses (10{,}000 resamples).\par}
\end{table}

\subsubsection{Reading the table}
Three conclusions follow, and we state them plainly.

\textbf{The 0.775 comes from copying templates, and we can now show it with numbers.} For all 150 test clauses the fine-tuned model outputs one of the 41 template sentences \emph{word for word}: 100\% of outputs are an exact copy of a template, the right one 70\% of the time and another category's template the other 30\%. The model never writes a sentence of its own. When it is wrong, it is wrong the way a classifier is wrong: a \emph{Price Restrictions} clause about a 10\% rate cap is answered with the \emph{ROFR} template. Fine-tuning did not teach the model to summarize; it taught it to pick one of 41 sentences.

\textbf{The last row of Table~\ref{tab:sumeval} settles it.} The plain category template, with no model at all, scores 0.115 against the clause-specific references, the same as the fine-tuned model's 0.116 within the margin of error. A simple lookup table gets the model's whole score. Whatever the fine-tuned summarizer learned adds nothing measurable beyond ``print this category's sentence''.

\textbf{Fine-tuning made the model worse at the real task.} Against the clause-specific references, the \emph{original} model (before fine-tuning) scores 0.299, more than twice the fine-tuned model's 0.116. This does not mean the original FLAN-T5 is a good legal summarizer: it is not. As Table~\ref{tab:sumqual} shows, it mostly repeats the source text, and repeating the source raises a word-overlap score against a reference that uses the same words as the clause. The honest conclusion is narrower but still uncomfortable: our fine-tuning removed whatever clause-specific content the original model had, and replaced it with smooth, general sentences that score well only against references made of the same sentences.

\begin{table}[!htbp]
\centering
\caption{Fine-tuned vs.\ zero-shot output}
\label{tab:sumqual}
\small
\begin{adjustbox}{max width=\textwidth}
\begin{tabular}{>{\raggedright\arraybackslash}p{3.2cm} >{\raggedright\arraybackslash}p{10.3cm}}
\toprule
\textbf{Input clause} & ``Any attempted assignment or delegation in violation of this section by either party without the prior written consent of the other will be void.''\\
\midrule
Gold template & You cannot transfer this contract to someone else without permission.\\
Fine-tuned    & You cannot transfer this contract to someone else without permission. \emph{(exact template)}\\
Zero-shot     & ``a) Any attempted assignment or delegation in violation of this section \ldots will be void.'' \emph{(repeats the legal text)}\\
\midrule
Clause-specific reference & Neither side may hand its rights or duties to a third party without the other's written consent, which cannot be unreasonably refused.\\
\bottomrule
\end{tabular}
\end{adjustbox}
\end{table}

\subsubsection{What we therefore claim}
What we can claim is a working \textbf{``sort into a category, then show a template''} system: it finds a clause's category and shows a plain-English sentence for that category, written in advance by people. That is still useful in an app: a signer who reads ``You cannot transfer this contract to someone else without permission'' next to a block of legal text is better off than before. (Clausify shows this line on every card, labelled as the summarizer's output and placed next to the fixed reason from the risk table; when the sentence belongs to another category, the card says so; Section~\ref{sec:clausify}.) We do \textbf{not} claim free-text legal summarizing, and the 0.775 should never be used as evidence of it. The number that describes how well the model summarizes individual clauses is 0.116, and the honest conclusion is that this part is a proof of concept, not a finished result.

Moving past this needs target sentences that differ from clause to clause, which is exactly what the trial set above gives a small sample of: the 150 references took a few hours to draft. Section~\ref{sec:futurework} describes the larger version.

\section{Analysis of Design Solutions}
This section looks at the chosen design as a whole: how many clauses get through the whole system, whether its mistakes happen where they would hurt a signer, what the mistakes look like, and how the ensemble behaves on rare categories.

\subsection{End-to-End: What Actually Reaches the User}
\label{sec:e2e}

Every number so far measured one step on its own, with the other steps out of the way. Presence was scored on its own; extraction was scored with a window already known to contain the answer; summarizing was scored on a real clause. That is the right way to test each part separately, but none of these tells us what the whole system does to a contract. The steps run one after another, and errors add up: a category the presence step misses is a clause the span step never sees.

So we ran the whole system on all 102 test contracts and asked one question about each of the 1{,}348 real clauses: does it get through \emph{all three} steps? A clause counts only if (i) the presence step finds its category, (ii) the span step, reading the window the presence step chose (never a window chosen using the real answer), returns text that overlaps the real clause with token-F1 $\geq$ 0.5, and (iii) the risk level given is the correct one. Table~\ref{tab:e2e} shows how many clauses are lost at each step.

\begin{table}[!htbp]
\centering
\caption{End-to-end clause survival}
\label{tab:e2e}
\small
\begin{adjustbox}{max width=\textwidth}
\begin{tabular}{llcc}
\toprule
\textbf{Step} & \textbf{Condition} & \textbf{Clauses left} & \textbf{Share of all}\\
\midrule
 -- & real clauses in the test set & 1{,}348 & 100.0\%\\
2A Presence & category found & 1{,}057 & 78.4\%\\
2B Span & extracted text overlaps the real clause at F1 $\geq$ 0.5 & 296 & 22.0\%\\
Risk & correct risk level given & 296 & 22.0\%\\
\midrule
\multicolumn{2}{l}{\textbf{Clauses that get through every step}} & \textbf{296} & \textbf{22.0\%}\\
\multicolumn{2}{l}{\emph{95\% CI, bootstrap over contracts}} & & $[19.8\%,\ 24.2\%]$\\
\bottomrule
\end{tabular}
\end{adjustbox}

{\footnotesize\singlespacing \emph{Note.} AND-ensemble presence rule. The last column is the share of real clauses still left after that step.\par}
\end{table}

\textbf{22.0\% is the number that describes what a user actually gets}, and it is far below what the separate scores suggest. Multiplying the separate figures, 78.4\% presence recall times 82.8\% span overlap, predicts about 65\%. The system delivers about a third of that. This gap is the reason for running the test, so it needs to be explained, not just reported.

The risk step is the one place where nothing is lost: the level is a fixed lookup for each category, so it is right exactly when the category is right, and it adds no errors of its own. That comes from the design, not from good performance, and ``correct'' here only means ``matches our risk table'', which, as Section~\ref{sec:taxvalid} notes, has not been checked by an expert.

\subsubsection{Where the clauses are lost}
The big loss happens at the span step, and it has two separate causes. We separated them by checking, for every clause whose category was found, whether the window chosen by the presence model really contained the clause.

\begin{table}[!htbp]
\centering
\caption{Span-stage loss breakdown}
\label{tab:e2eattr}
\small
\begin{adjustbox}{max width=\textwidth}
\begin{tabular}{lcc}
\toprule
 & \textbf{Count} & \textbf{Share}\\
\midrule
Chosen window \emph{contains} the real clause & 730 & 69.1\%\\
\quad and the span model found it (F1 $\geq$ 0.5) & 289 & 39.6\% of 730\\
\quad but the span model still missed it & 441 & 60.4\% of 730\\
Chosen window does \emph{not} contain the real clause & 327 & 30.9\%\\
\quad but the span still scored $\geq$ 0.5 (same wording elsewhere) & 7 & \\
\bottomrule
\end{tabular}
\end{adjustbox}
\end{table}

\textbf{Cause one: the window with the highest score is often the wrong window.} We chose the maximum rule (Section~\ref{sec:mil}) because it removes the search limit (if any window fires, the category is reported), and Section~\ref{sec:mil} discussed what that costs in \emph{precision}. What we did not expect is that the highest-scoring window is also being used to \emph{locate} the clause, and it is poor at that: 30.9\% of the time it does not contain the clause at all. The model was never trained for this. It was trained to answer ``does this window contain a clause of this category?'', and a window that talks about liability in general can score higher than the one that actually limits it. For those 327 clauses the span step is given the wrong 2{,}000 characters, and no extraction model could recover them.

\textbf{Cause two: the span model is much weaker on windows it was not trained on.} Even when the window does contain the clause, extraction succeeds only 39.6\% of the time, compared with 82.8\% in the separate test of Section~\ref{sec:span}. The reason is the way focus windows were built (Section~\ref{sec:spanmethod}). In training, every example was centred so the answer started about 150 characters into a 1{,}200-character window. That is allowed: it uses the answer to build a training example, never as an input when the system runs, and Section~\ref{sec:leakage} explains why. But allowed is not the same as harmless: the model has only ever seen windows where the answer is near the start, while in the real system it reads plain 1{,}200-character chunks where the answer can be anywhere. The average token-F1 falls from 0.764 to 0.383. The centring did not leak the test set, but it did quietly narrow the kind of text the model can handle.

\subsubsection{What this changes}
Three things follow, and we prefer to say them clearly rather than let the separate scores stand alone.

First, \textbf{the span scores in Section~\ref{sec:span} are an upper limit on real extraction quality, not a description of it}. They answer ``how good is the extractor when given the right window?'', which is a fair question about the model but a misleading one about the system.

Second, \textbf{the most useful next step is not a bigger model}. Cause two is a data problem with a cheap fix (train the span model on windows where the answer is at different positions, or on the same plain chunks it sees in the real system), and cause one calls for choosing the extraction window with a model trained for that job instead of reusing the presence score. Both fit within this project's computing budget; neither needs more parameters.

Third, \textbf{Clausify's output should be read as a shortlist, not as exact extraction}. When the system shows a clause card, the category is right about 78\% of the time the clause was present, but in this test the span model's quoted text was a good copy of the real clause only about one time in five overall. The app has since changed how it chooses the quote: it now quotes the whole paragraph that the presence model scores highest, which fixed the clear failures we saw in our own tests, but this new method has not been measured on the test set. The app tells users that it is not legal advice and that the quoted text may not be the exact clause the badge refers to (Section~\ref{sec:demolimits}).
\subsection{Risk-Weighted Evaluation: Does the Metric Match the Purpose?}
\label{sec:riskeval}

Micro-F1 over 4{,}182 decisions treats every decision as equally important. A person signing a contract does not. The whole idea of the risk layer (Section~\ref{sec:risk}) is that some categories matter more than others, so the evaluation should ask whether the system's mistakes fall where they do the least harm. We had not asked this before, and the answer turns out to reverse the recommendation that the overall score gives.

Table~\ref{tab:riskeval} splits the 41 categories into the three risk levels of our own risk table and scores each model within each level.

\begin{table}[!htbp]
\centering
\caption{Results by risk band}
\label{tab:riskeval}
\small
\begin{adjustbox}{max width=\textwidth}
\begin{tabular}{llccccc}
\toprule
\textbf{Band} & \textbf{Model} & \textbf{Precision} & \textbf{Recall} & \textbf{F1} & \textbf{Found} & \textbf{Missed}\\
\midrule
\multirow{4}{*}{\shortstack[l]{\textbf{High}\\8 cats\\176 positives}}
 & Transformer  & 0.392 & \textbf{0.841} & 0.534 & 148 & \textbf{28}\\
 & TF--IDF      & 0.540 & 0.648 & 0.589 & 114 & 62\\
 & \textbf{AND ensemble} & 0.581 & 0.631 & 0.605 & 111 & \textbf{65}\\
 & OR ensemble  & 0.379 & \textbf{0.858} & 0.526 & 151 & \textbf{25}\\
\addlinespace
\multirow{4}{*}{\shortstack[l]{\textbf{Medium}\\22 cats\\498 positives}}
 & Transformer  & 0.429 & 0.886 & 0.578 & 441 & 57\\
 & TF--IDF      & 0.621 & 0.723 & 0.668 & 360 & 138\\
 & \textbf{AND ensemble} & 0.655 & 0.679 & 0.667 & 338 & 160\\
 & OR ensemble  & 0.424 & 0.930 & 0.582 & 463 & 35\\
\addlinespace
\multirow{4}{*}{\shortstack[l]{\textbf{Low}\\11 cats\\674 positives}}
 & Transformer  & 0.800 & 0.953 & 0.870 & 642 & 32\\
 & TF--IDF      & 0.909 & 0.917 & 0.913 & 618 & 56\\
 & \textbf{AND ensemble} & 0.923 & 0.902 & 0.912 & 608 & 66\\
 & OR ensemble  & 0.792 & 0.967 & 0.871 & 652 & 22\\
\bottomrule
\end{tabular}
\end{adjustbox}

{\footnotesize\singlespacing \emph{Note.} Threshold 0.5. ``Missed'' counts real clauses of that level that the model did not report.\par}
\end{table}

\subsubsection{The headline configuration is the worst one for High-risk clauses}
Look at the High level. The AND-ensemble (the setup we report as our main result, because of its overall micro-F1) \textbf{misses 36.9\% of High-risk clauses}. The transformer alone, which the previous sections showed to be the weaker model, misses only \textbf{15.9\%}. The more permissive OR rule, which comes \emph{last} on overall micro-F1 (0.692, tied with the transformer alone), misses only \textbf{14.2\%}.

In numbers, out of 176 High-risk clauses in 102 contracts: the ensemble finds 111 and misses 65; OR finds 151 and misses 25. Choosing the ensemble over OR because of micro-F1 costs a signer 40 High-risk clauses (unlimited liabilities, IP assignments, non-competes), in exchange for a leaderboard position that Section~\ref{sec:bootstrap} already showed is a statistical tie.

The reason is the one found in Section~\ref{sec:raretail}, and it adds up here. Several High-risk categories are uncommon: \emph{Most Favored Nation} is among the ten rarest categories, and \emph{Liquidated Damages} (11 test examples) and \emph{Irrevocable or Perpetual License} (17) are also below average. The AND rule is hardest to satisfy exactly where there is the least training data, so the risk level that needs recall the most is the one where requiring agreement hurts it most. The Low level shows the opposite: 674 of the 1{,}348 real clauses are Low-risk (dates, party names, governing law), and every model scores about 0.9 there. \textbf{The overall micro-F1 is dominated by the clauses that matter least.}

\subsubsection{What we conclude, and what we do not}
We do not change the reported setup because of this table, and we should say why: choosing OR \emph{after} seeing which rule does best on the High level of the test set would be choosing a setup using the test set, exactly what Section~\ref{sec:valprotocol} says this project avoided. The result stands as a diagnosis, not as a new tuning.

What it does show is that \textbf{our evaluation score and our stated goal never matched}, and that the mismatch is big enough to reverse the ranking of the models. A system meant to warn a signer before signing should be chosen by how many High-risk clauses it finds, or by a risk-weighted F1, not by an average over 41 categories where half of the real clauses are Low-risk ones such as dates and party names. Section~\ref{sec:futurework} describes the proper version of this test: decide the risk-weighted goal first, then choose the rule on a separate validation split, then report on the test set once.

Finally, this is one more reason to take Section~\ref{sec:taxvalid} seriously. Every number in Table~\ref{tab:riskeval} is computed using \emph{our own} choice of which categories are High, Medium and Low. If a lawyer disagrees about which categories are High, the table changes too. The risk-weighted test inherits all the warnings about the unchecked risk table and adds no independent evidence for it.
\subsection{What the errors actually look like}
\label{sec:errorexamples}
The confusion matrix counts 309 false positives and 291 false negatives, but counts alone do not show whether the model fails for silly or understandable reasons. Table~\ref{tab:errors} shows real cases from the test contracts: false positives sorted by the model's own confidence, because a system's most confident mistakes are the most revealing.

\begin{table}[!htbp]
\centering
\caption{Representative errors}
\label{tab:errors}
\scriptsize
\renewcommand{\arraystretch}{1.25}
\begin{adjustbox}{max width=\textwidth}
\begin{tabular}{@{}p{2.1cm} >{\raggedright\arraybackslash}p{1.5cm} >{\raggedright\arraybackslash}p{9.6cm}@{}}
\toprule
\textbf{Category} & \textbf{Scores} & \textbf{Evidence}\\
\midrule
\multicolumn{3}{@{}l}{\textbf{False positives}: reported, but the category is not labelled in this contract}\\
\midrule
Insurance & 0.82 / 0.96 & ``\ldots a fixed amount of [*] of gross sales to cover freight, postage, \textbf{insurance costs} on shipments to such Third Party, packing costs\ldots'' \emph{A royalty clause that lists insurance as a cost. The word is there; the duty to insure is not.}\\
\addlinespace
Governing Law & 0.81 / 0.95 & ``\ldots there shall arise the need to decide the question by what municipal or national law this Agreement \ldots is \textbf{governed}, the following facts shall be excluded\ldots'' \emph{A clause about how to work out the governing law, not one that states it. A borderline case for the labels.}\\
\addlinespace
Parties & 1.00 / 0.98 & An addendum whose first window looks like the top of a contract. \emph{Both models react to the layout; the labels put Parties somewhere else.}\\
\midrule
\multicolumn{3}{@{}l}{\textbf{False negatives}: labelled as present, but both models scored it low}\\
\midrule
Liquidated Damages & 0.21 / 0.02 & ``\ldots agrees to pay a \textbf{cancellation fee} of Five Thousand Dollars (\$5{,}000) per month multiplied by the number of months remaining\ldots'' \emph{A clear liquidated-damages term that never uses the phrase. The word model has nothing to match; the transformer, with 50 training contracts for this category, has not learned the idea.}\\
\addlinespace
Most Favored Nation & 0.22 / 0.02 & ``\ldots Sponsor shall have the right to approve or disapprove any additional sponsor\ldots unless another\ldots'' \emph{In meaning an MFN-style protection, but written without any of the category's usual words. This category has 21 training contracts and scores 0.000 overall.}\\
\addlinespace
Third Party Beneficiary & 0.26 / 0.02 & ``Notwithstanding the foregoing, UTK may assign this Agreement or any portion of its Compensation \ldots to its subsidiaries.'' \emph{This reads like an assignment clause, not a third-party-beneficiary clause. Some rare-category ``failures'' are disagreements with the labels, not model errors.}\\
\midrule
\multicolumn{3}{@{}l}{\textbf{Span errors}: category found correctly, but the extracted text is wrong (Section~\ref{sec:e2e})}\\
\midrule
Minimum Commitment & F1 0.00 & The window \emph{does} contain the clause, but the model returns only ``\texttt{the}''. \emph{The start and end of the answer fall on the same word, and nothing in the decoding stops this.}\\
\addlinespace
Effective Date & F1 0.06 & The window does \emph{not} contain the real clause, yet the model confidently returns ``7th day of september, 1999''. \emph{When Stage 2A gives the wrong window, Stage 2B still answers: it has no way to say ``not here''.}\\
\addlinespace
Insurance & F1 0.20 & The real clause is the insurance duty; the model returns the indemnity clause next to it. \emph{Close in topic, different in law.}\\
\bottomrule
\end{tabular}
\end{adjustbox}

{\footnotesize\singlespacing \emph{Note.} Text shortened. Scores are the TF--IDF and transformer probabilities.\par}
\end{table}

Three patterns stand out. \textbf{Matching words without the actual duty} explains many false positives: a contract that mentions insurance costs is not a contract with an insurance requirement, and a bag-of-words model reading the whole document cannot tell the difference. \textbf{The idea without the usual words} explains many misses: a \$5{,}000-per-month cancellation fee \emph{is} liquidated damages, but a model that has seen only about fifty contracts with this category, most of them using the usual words, will not learn that. And \textbf{disagreement with the labels} is a real part of the rare-category problem: the \emph{Third Party Beneficiary} case above looks like a labelling mistake rather than a model miss, so the 0.000 F1 scores in Section~\ref{sec:raretail} are partly model failure and partly limits of the labels. We did not try to measure how much of each; that would need the contracts to be labelled again, which we are not qualified to do.

The span errors point to something easier to fix. The near-empty answer ``\texttt{the}'' and the confident answer from a window that cannot contain the clause have the same cause: the extraction model has no way to say \emph{there is nothing here}. Training with unanswerable examples, as in SQuAD~v2, would give it exactly that, and it is listed as future work.
\subsection{Does the AND rule punish rare categories harder than common ones?}
\label{sec:raretail}
Section~\ref{sec:bootstrap} supports the AND-ensemble because of the kind of mistakes it makes: it turns the transformer's many false alarms into balanced precision and recall. But that argument was made on the overall numbers, and an overall number can hide a trade that is not shared evenly. It is worth asking where the AND rule's extra precision comes from, because there is a clear reason to worry about the rare categories. Requiring two models to agree is hardest exactly where each model has the least training data, so the rarest categories are where the second model is least likely to agree, and agreement is what the rule requires.

We sorted the 41 categories by how often they appear in the \emph{training} data, took the ten rarest (11 to 42 positive training contracts, 70 positive test decisions in total), and compared the AND rule with the transformer alone on these ten and on the other 31 categories separately. Table~\ref{tab:raretail} gives the result, and the two halves tell different stories.

\begin{table}[!htbp]
\centering
\caption{Rare vs.\ common categories}
\label{tab:raretail}
\small
\begin{adjustbox}{max width=\textwidth}
\begin{tabular}{llccccc}
\toprule
\textbf{Group} & \textbf{Model} & \textbf{Precision} & \textbf{Recall} & \textbf{F1} & \textbf{TP} & \textbf{FN}\\
\midrule
\multirow{3}{*}{Rarest 10 (70 test positives)}
 & Transformer & 0.275 & 0.557 & \textbf{0.368} & 39 & 31\\
 & TF--IDF     & 0.429 & 0.300 & 0.353 & 21 & 49\\
 & \textbf{AND ensemble} & 0.621 & 0.257 & 0.364 & 18 & 52\\
\addlinespace
\multirow{3}{*}{Other 31 (1{,}278 test positives)}
 & Transformer & 0.577 & 0.933 & 0.713 & 1{,}192 & 86\\
 & TF--IDF     & 0.753 & 0.838 & 0.793 & 1{,}071 & 207\\
 & \textbf{AND ensemble} & 0.777 & 0.813 & \textbf{0.795} & 1{,}039 & 239\\
\bottomrule
\end{tabular}
\end{adjustbox}

{\footnotesize\singlespacing \emph{Note.} Micro-averaged within each group. Rarest 10 = fewest positive training contracts (11 to 42).\par}
\end{table}

\textbf{On the common categories the ensemble works as intended.} Compared with the transformer alone it gains 20.0 points of precision and loses 12.0 points of recall, and F1 rises from 0.713 to 0.795. That is a good trade, and it is where all of the ensemble's overall advantage comes from.

\textbf{On the rare categories it gains nothing: it is a straight swap.} The AND rule gains 34.6 points of precision by giving up 30.0 points of recall, and its F1 ends at 0.364, slightly below the transformer alone (0.368). Every point of precision on the rare categories is paid for with recall, with nothing left over. In numbers, the transformer alone finds 39 of the 70 rare clauses; requiring agreement cuts this to 18. The TF--IDF model alone is now the weakest of the three on these categories (F1 0.353), and even so, its vetoes decide what the ensemble reports there.

The details per category are clearer than the average. On five of the ten rarest categories the AND rule scores exactly 0.000, and on three of those (\emph{Price Restrictions}, \emph{Unlimited/All-You-Can-Eat License}, \emph{No-Solicit of Customers}) the transformer alone was finding something (F1 0.087, 0.500, 0.323) before the AND rule removed it. So the rule does not just fail to help on rare categories; on these categories it removes clauses that one of the models had found.

\textbf{Why this matters for real use, not just for the table.} The overall micro-F1 counts all 4{,}182 decisions equally, but a signer does not. Missing a rare clause is not clearly less costly than missing a common one, and it may be much worse: \emph{Most Favored Nation} is among the ten rarest categories and is one of our eight High-risk categories, and both models score 0.000 on it. A system whose rule favours precision on common categories, even a rule with no settings, is improving the score rather than reducing harm. We report the ensemble as our main setup because that is what the overall test chooses. But on this evidence, a system meant to protect signers would do better to \emph{use the more permissive rule for High-risk categories and the AND rule for the rest}, accepting more false alarms exactly where a missed clause costs the most. We have not built that mixed rule: building it and testing it properly is future work (Section~\ref{sec:futurework}), not something to add after looking at the test set.

\section{Final Design Adjustments}
Two design choices made before training, max-pooling as the aggregation rule and 0.5 as the decision threshold, were never tuned. After the main evaluation both were revisited, to determine whether adjusting them would change the conclusions or the deployed system.

\subsection{Ablation: was max-pooling the right aggregation rule?}
\label{sec:pooling}
Section~\ref{sec:mil} chose max-pooling because it \emph{is} the multiple-instance decision rule \cite{maron1998mil}, a bag is positive if any instance is, and Sections~\ref{sec:mil} and~\ref{sec:budget} then blamed max-pooling for the transformer's inflated false positives. Both statements were made without ever trying an alternative, which is not good enough for a claim the project leans on. We therefore scored every window once and compared aggregation functions on \emph{identical} window scores, so the only thing varying is the rule.

\begin{table}[!htbp]
\centering
\caption{Aggregation ablation}
\label{tab:pooling}
\small
\begin{adjustbox}{max width=\textwidth}
\begin{tabular}{lccccc}
\toprule
\textbf{Aggregation} & \textbf{F1@0.5} & \textbf{P@0.5} & \textbf{R@0.5} & \textbf{Best F1} & \textbf{at threshold}\\
\midrule
Max (used in this project) & 0.6943 & 0.562 & 0.908 & \textbf{0.7504} & 0.90\\
Top-3 mean                & \textbf{0.7144} & 0.667 & 0.769 & 0.7210 & 0.40\\
Top-5 mean                & 0.6661 & 0.719 & 0.620 & 0.7192 & 0.30\\
Mean over all windows     & 0.1664 & 0.946 & 0.091 & 0.7114 & 0.10\\
Noisy-OR                  & 0.5897 & 0.423 & 0.973 & 0.7022 & 0.95\\
Second-highest window     & 0.6832 & 0.662 & 0.705 & 0.7016 & 0.30\\
\midrule
\emph{TF--IDF baseline (no windows)} & \textbf{0.7747} & 0.742 & 0.810 & -- & -- \\
\bottomrule
\end{tabular}
\end{adjustbox}

{\footnotesize\singlespacing \emph{Note.} ``Best'' is each rule's micro-F1 at its own optimal threshold, found by sweeping the test set: a diagnostic upper bound, not a configuration we adopt.\par}
\end{table}

Three things come out of this, and one of them corrects an earlier claim of ours.

\textbf{At the threshold we actually used, max-pooling was not the best choice.} Top-3 mean scores 0.7144 against max's 0.6943 at the 0.5 boundary: a fifth of the gap to the baseline, available for free from a one-line change. Averaging the three highest-scoring windows is a natural softening of the MIL rule: it still asks whether evidence is concentrated somewhere in the document, but requires that evidence to appear in more than one place, which is exactly the guard against a single spurious window that Section~\ref{sec:mil} identified as max-pooling's weakness. Mean-over-all-windows collapses at 0.5 (recall 0.091) for the obvious reason, an average over 35 windows of which one is positive cannot reach 0.5, and its best threshold is correspondingly 0.10.

\textbf{Once each rule is given its own threshold, max-pooling is the best of the six.} At 0.90 it reaches 0.7504, ahead of top-3 mean's 0.7210. So max-pooling is not a poor aggregation rule; it is a rule whose output distribution is badly matched to a 0.5 decision boundary, which is the same finding Section~\ref{sec:threshold} reached from the other direction. The two results are one result: the transformer's reported deficit is substantially an operating-point problem, not an aggregation problem.

\textbf{But no aggregation rule reaches the baseline.} The best cell in the entire table (max-pooling at a threshold tuned on the test set, which is already a generous upper bound) is 0.7504 against the baseline's 0.7747. Every alternative rule, at every threshold we swept, remains behind. This matters more than the first two points, because it closes off the last available explanation for RQ1: the windowed transformer does not lose because of its training budget (Section~\ref{sec:budget}), or its seed (Section~\ref{sec:seeds}), or its aggregation rule, or its threshold. It loses.

\textbf{A correction to Section~\ref{sec:budget}.} There we wrote that ``the ceiling here is the aggregation rule, not the encoder underneath it'', inferring it from the precision/recall shift when training data was doubled. Table~\ref{tab:pooling} shows that inference was wrong: changing the aggregation rule (through six variants, each at its best threshold) moves micro-F1 by at most 0.056 and never past the baseline. The aggregation rule is a real but secondary factor. We have left the original sentence in place with this correction attached, on the same principle as elsewhere in this chapter: an inference we drew from indirect evidence and later measured properly is more useful to a reader than a silently amended one.

A caveat on how far this ablation goes. It varies only the function applied to window scores. It does not vary the window geometry (2{,}000 characters, 1{,}500 stride), which would require re-windowing and retraining, nor does it try a learned aggregator such as attention pooling over window representations. Both remain untested, and attention pooling in particular is the obvious next thing to try, since it could learn to weight windows by relevance instead of taking a fixed statistic.
\subsection{The Threshold and What the Confidence Score Means}
\label{sec:threshold}

Every presence decision in this project, and every clause card in Clausify, comes from comparing a score with 0.5. Section~\ref{sec:valprotocol} notes that this threshold was never chosen: it is simply the default for a two-class model. This raises two questions we had not asked: is 0.5 a good place to draw the line, and does the score being compared with it mean anything? Figure~\ref{fig:threshcal} answers both.

\begin{figure}[!htbp]
\centering
\includegraphics[width=\textwidth]{threshold_calibration.png}
\caption{Threshold and calibration}
\label{fig:threshcal}
\end{figure}

\subsubsection{0.5 is a poor operating point for the transformer}
Trying other thresholds after the fact shows that the transformer reaches a micro-F1 of \textbf{0.753 at a threshold of 0.9}, against \textbf{0.692} at the default: a gain of 0.061, which closes most of its gap with the TF--IDF baseline (0.775). The reason is the one Section~\ref{sec:mil} predicted: the maximum over about 35 windows is pushed upwards by design, so the scores bunch up near 1 and a line drawn at 0.5 lets far too much through. Precision is 0.558 at 0.5 and 0.726 at 0.9.

\textbf{We are not switching to 0.9.} It was found by looking at the test set, and using it would turn a fair held-out test into a tuned one, which is exactly what Section~\ref{sec:valprotocol} says this project avoided. The honest statement is narrower and more useful: \emph{a real part of the transformer's gap comes from a threshold nobody chose, not from an inability to tell clauses apart}. So the comparison in Section~\ref{sec:bootstrap} sets a baseline whose scores already suit 0.5 against a transformer whose scores do not. A fair rematch would choose each model's threshold on a validation split taken from the 408 training contracts, and we did not do that. Section~\ref{sec:budget} showed the gap is not caused by too little training data; this section shows part of it may be caused by the threshold, and that warning belongs next to the finding.

The AND-ensemble behaves the other way: 0.5 is already its best threshold (0.779), and raising it makes things much worse. At 0.7 micro-F1 falls to 0.611 and High-risk recall to 0.165, because asking \emph{both} models to pass a high bar is close to asking both to be very sure at once. The dashed curves in Figure~\ref{fig:threshcal} show the conflict: for the transformer, the threshold with the best micro-F1 (0.9) lowers High-risk recall from 0.841 to 0.716. No single threshold is best both for the overall score and for the clauses that matter most.

\subsubsection{The confidence bar in the application is not a probability}
Clausify shows the presence score to the user as a confidence bar. A reliability diagram tests whether that bar means what a user would think it means, and it does not.

The transformer's score has an \textbf{Expected Calibration Error of 0.205} and a Brier score of 0.195. It is too confident everywhere. Among predictions between 0.5 and 0.6, the clause is really there only \textbf{16.5\%} of the time, and even in the top group (scores above 0.9, which the app shows as an almost full bar) it is there \textbf{72.6\%} of the time. A user who reads ``96\% confident'' is really looking at about a 73\% chance. Again the cause is max-pooling: the maximum over dozens of windows is not a probability, and nothing in training asked it to be one.

The TF--IDF baseline is also off (ECE 0.172), but mostly in the safer direction: it is \emph{less} confident than it should be in the middle range. Its predictions between 0.7 and 0.8 are right 89.6\% of the time, and those between 0.8 and 0.9 are right 97.8\% of the time.

The practical effect on the app is direct, and we record it in Section~\ref{sec:demolimits}: the score shown on every card is not a probability of being right. Fixing this needs no retraining. Platt scaling or isotonic regression, fitted on a validation split, would turn the scores into real probabilities, and it is the cheapest single improvement to the interface. We had no validation split set aside to fit it on, which follows from the gap in Section~\ref{sec:valprotocol} rather than being a separate mistake.

\section{Statistical Analysis}
Our main statistical tool is a paired bootstrap over contracts \cite{efron1979bootstrap}: we resample the 102 test contracts with replacement 10{,}000 times and score every model on the same resample, so the 41 linked decisions from each contract stay together. We then measure three sources of variation the bootstrap cannot see: the amount of training data, the random seed and the train/test split.

\subsection{How much of that ranking is real?}
\label{sec:bootstrap}
Table~\ref{tab:ensemble} is a ranking, and a ranking tempts the reader to treat the order itself as a result. On 102 test contracts, a lead of 0.779 over 0.775, about four decisions in a thousand, is too small to read that way without checking, so we measured it.

The natural tool is a \textbf{bootstrap confidence interval}, but what we resample matters. The 4{,}182 test decisions are not independent: 41 of them come from the same contract, and an unusual contract (very long, or oddly written) moves all 41 together. Resampling single (contract, category) decisions would ignore this and give intervals that are far too narrow. So we resample whole \emph{contracts} with replacement (a cluster bootstrap): we draw 102 contracts from the 102 we have, recompute micro-F1 for both models, and repeat 10{,}000 times. Both models are scored on exactly the same resample each time, so the difference is \emph{paired}, which makes the interval for the difference much narrower than the two separate intervals would suggest.

\begin{table}[!htbp]
\centering
\caption{Cluster bootstrap results}
\label{tab:bootstrap}
\small
\begin{adjustbox}{max width=\textwidth}
\begin{tabular}{lccc}
\toprule
\textbf{Comparison} & \textbf{Difference in micro-F1} & \textbf{95\% CI} & \textbf{$P(>0)$}\\
\midrule
AND-ensemble vs.\ TF--IDF baseline & $+0.0042$ & $[-0.0037,\ +0.0115]$ & 0.855\\
TF--IDF baseline vs.\ transformer  & $+0.0824$ & $[+0.0647,\ +0.1013]$ & 1.000\\
AND-ensemble vs.\ transformer      & $+0.0866$ & $[+0.0705,\ +0.1035]$ & 1.000\\
\bottomrule
\end{tabular}
\end{adjustbox}

{\footnotesize\singlespacing \emph{Note.} 10{,}000 resamples. $P(>0)$ is the share of resamples in which the first model wins.\par}
\end{table}

The three comparisons do not come out the same way, and that difference is the point.

\textbf{The ensemble's win over the baseline is not proven.} The interval is $[-0.0037, +0.0115]$, which includes zero, and the ensemble is ahead in 85.5\% of resamples: more often than not, but well short of the 97.5\% a clear win would need. On their own the two models score 0.779 (95\% CI $[0.762, 0.795]$) and 0.775 (95\% CI $[0.758, 0.791]$), and these intervals almost completely overlap. On this test set we cannot tell the AND-ensemble and the TF--IDF baseline apart, and we state this as a measured fact, not as a hedge. The top two rows of Table~\ref{tab:ensemble} should be read as a tie.

\textbf{The transformer's loss, on the other hand, is real.} The baseline beats the window-based transformer by 0.0824, with an interval of $[+0.0647, +0.1013]$ that is nowhere near zero, and it wins in all 10{,}000 resamples. This matters because the project's two main claims do not stand on equal ground. The finding we trust most, that a whole-document word model beats a window-based DistilBERT at deciding which clauses a contract contains, passes the test easily. The finding we were already careful about, that the AND-ensemble is ahead of the baseline, does not pass. It is good that our caution pointed the right way, but it needed to be checked, not assumed.

None of this weakens the case for the ensemble, because that case never rested on 0.004 of micro-F1. The AND rule earns its place by the \emph{kind of mistakes} it makes, not by its rank: it turns the transformer's many false alarms into balanced precision and recall (Table~\ref{tab:metrics}). Sections~\ref{sec:seeds} and~\ref{sec:raretail} look at what this costs.
\subsection{How much of this moves if we only change the random seed?}
\label{sec:seeds}
Section~\ref{sec:bootstrap} measured one kind of uncertainty: how much the score depends on \emph{which contracts} ended up in the test set. There is a second kind it cannot see. Every result so far comes from one training run per model, and training has randomness in it: the starting weights of the classifier, the order of the data and the sample of negative windows all depend on the random seed. If changing only the seed moves the score as much as the effects we are discussing, then a ranking built on single runs is reporting noise.

This is easy to check, so we trained the presence transformer three times with 24{,}000 training examples, changing nothing but the seed (42, 43 and 44), and tested each run on the same 102 contracts. Seed 42 is the October 2026 run used everywhere else in this chapter; seeds 43 and 44 were trained in August 2026 with the same settings. Table~\ref{tab:seeds} shows the spread.

\begin{table}[!htbp]
\centering
\caption{Seed variance}
\label{tab:seeds}
\small
\begin{adjustbox}{max width=\textwidth}
\begin{tabular}{lcccccc}
\toprule
\textbf{Model} & \textbf{Seed 42} & \textbf{Seed 43} & \textbf{Seed 44} & \textbf{Mean} & \textbf{SD} & \textbf{Range}\\
\midrule
Transformer (max-pool) & 0.6924 & 0.6928 & 0.7082 & 0.6978 & 0.0090 & 0.0158\\
AND ensemble           & 0.7789 & 0.7774 & 0.7803 & 0.7789 & 0.0015 & 0.0029\\
\midrule
TF--IDF baseline       & \multicolumn{3}{c}{0.7747 (deterministic)} & 0.7747 & -- & -- \\
\bottomrule
\end{tabular}
\end{adjustbox}

{\footnotesize\singlespacing \emph{Note.} All three runs use 24{,}000 training examples and differ only in seed. Seed 42 is the October 2026 run; seeds 43 and 44 are August 2026 runs. The TF--IDF baseline has no randomness, so it appears once.\par}
\end{table}

\textbf{Seed noise alone does not explain the ensemble's lead, but other noise does.} The AND-ensemble's lead over the baseline is $+0.0042$. Across seeds, the ensemble's standard deviation is only $0.0015$ and its range is $0.0029$, so the lead is larger than the seed-to-seed spread, and the ensemble beats the baseline with every seed (by up to $+0.0056$ with seed 44). That is a real change from the August models, where the lead ($+0.0014$) was smaller than the seed noise. But seed noise is the smallest of the uncertainties we measured. The bootstrap interval still includes zero (Section~\ref{sec:bootstrap}), and a different train/test split moves the baseline by $0.0063$ (Section~\ref{sec:splitvar}), 1.5 times the lead. So we still read the top two rows as a tie.

\textbf{The finding we stand behind is not affected by seed noise.} The baseline beats the transformer by 0.0824, which is 9.2 times the transformer's seed standard deviation of 0.0090. Even the luckiest seed, at 0.7082, still loses to the baseline by 0.0665. No seed makes the window-based transformer competitive, and together with Section~\ref{sec:budget}, where more than doubling the training data also failed to close the gap, this leaves the design of the model as the explanation still standing after both alternatives were tested.

It is worth noticing which model is noisier, because it is not the one a reader might expect. The transformer alone varies about six times as much across seeds ($\mathrm{SD} = 0.0090$) as the ensemble that contains it ($0.0015$). Requiring agreement with a fixed word model calms the transformer's run-to-run changes as well as its false alarms: a false detection has to get past a second, fixed opinion, so most of the seed-driven changes in what the transformer reports are filtered out. This steadiness is a real property of the AND rule, although Section~\ref{sec:raretail} shows it costs recall on rare categories. It is a better argument for the ensemble than its rank ever was.

\textbf{What this check covers.} We re-ran only the presence transformer with new seeds. The span extractor and the summarizer were each trained twice with the \emph{same} seed (August and October 2026) and gave the same scores to within 0.001 (token-F1 0.763 and 0.764; ROUGE-L 0.775 both times). This limits run-to-run noise on this hardware but says nothing about seed variance, which remains unmeasured for those two models; we list it in Section~\ref{sec:threats} instead of assuming it is small.
\subsection{Was the transformer simply undertrained?}
\label{sec:budget}
There is one objection to Section~\ref{sec:bootstrap} that the bootstrap cannot answer, and it goes to the heart of the project's main claim. The TF--IDF baseline is trained on all 408 training contracts. The window-based transformer was not: to keep training possible on a laptop, it used a balanced sample of 24{,}000 of the 53{,}662 available training windows (Section~\ref{sec:prestrain}). So when the transformer loses by about 0.08 micro-F1, that comparison alone cannot tell us whether max-pooling over windows is really weaker for this task, or whether the transformer simply never saw 55\% of its own training data. The first is a finding about the method; the second is a side effect of our computing budget. They lead to very different conclusions.

The only way to settle this is to remove the difference, so in August 2026 we retrained the presence transformer on \textbf{all 53{,}662 training windows}, with every other setting the same: same windows, same negative sampling, same settings, same seed, same test. The run took 140.7 minutes on the same laptop, against about 50 minutes for the smaller version. Table~\ref{tab:budget} shows the result. Both runs in this table are from August 2026, so they are compared with each other, not with the October numbers.

\begin{table}[!htbp]
\centering
\caption{Training-budget control}
\label{tab:budget}
\small
\begin{adjustbox}{max width=\textwidth}
\begin{tabular}{lccccc}
\toprule
\textbf{Model} & \textbf{Training windows} & \textbf{micro-F1} & \textbf{Precision} & \textbf{Recall}\\
\midrule
Transformer (max-pool)  & 24{,}000 (44.7\%) & 0.6943 & 0.5620 & 0.9080\\
Transformer (max-pool)  & \textbf{53{,}662 (100\%)} & 0.6939 & 0.5546 & 0.9266\\
TF--IDF baseline        & all 408 contracts & 0.7747 & 0.7424 & 0.8101\\
\midrule
\multicolumn{5}{l}{\emph{Paired differences (95\% CI over 10{,}000 contract resamples)}}\\
\multicolumn{2}{l}{Full-data transformer $-$ subsampled transformer} & $-0.0004$ & \multicolumn{2}{l}{$[-0.0098,\ +0.0095]$}\\
\multicolumn{2}{l}{TF--IDF baseline $-$ full-data transformer} & $+0.0809$ & \multicolumn{2}{l}{$[+0.0649,\ +0.0980]$}\\
\bottomrule
\end{tabular}
\end{adjustbox}

{\footnotesize\singlespacing \emph{Note.} August 2026 runs, both tested in the same way on the 102 test contracts; differences use the paired bootstrap of Section~\ref{sec:bootstrap}. The October 2026 re-run of the 24{,}000-window model scores 0.6924.\par}
\end{table}

\textbf{Too little training data is not the explanation.} Using 124\% more training data changed micro-F1 by $-0.0004$, with an interval of $[-0.0098, +0.0095]$ that sits right across zero: the two budgets cannot be told apart. And the gap to the baseline did not move: $+0.0809$ with all the data against $+0.0805$ with the sample (August runs), with an interval nowhere near zero and the baseline winning in all 10{,}000 resamples. Whatever holds the window-based transformer back, it is not a shortage of training examples, and the finding this project relies on survives the check.

The \emph{shape} of the extra data's effect is worth a sentence, because it is not simply ``nothing happened''. Recall rose from 0.9080 to 0.9266 and precision fell from 0.5620 to 0.5546: more training data made the model fire more readily, not more accurately. That is exactly the failure mode Section~\ref{sec:mil} predicted for max-pooling: with a mean of 35 windows per contract, the document-level decision is a maximum over many chances to be wrong, and a better-trained window scorer that is still imperfect produces \emph{more} spurious maxima, not fewer. The ceiling here is the aggregation rule, not the encoder underneath it.\footnote{This inference turned out to be only partly right, and Section~\ref{sec:pooling} corrects it: swapping max-pooling for five alternatives moves micro-F1 by at most 0.056 and never past the baseline, so the aggregation rule is a secondary factor rather than the ceiling.}

One smaller effect is worth reporting even though it does not change our conclusion. The AND-ensemble built on the full-data transformer scores 0.7820, and its lead over the baseline is $+0.0073$ with a 95\% interval of $[-0.0000, +0.0140]$: it wins in 97.4\% of resamples, but the interval still touches zero. So more training data moves the ensemble's lead from ``clearly a tie'' to ``borderline''. We do not claim it: an interval whose lower end sits on zero does not pass the bar we set in Section~\ref{sec:bootstrap}, and we will not lower that bar because the number moved our way. What we can say is that the ensemble's balanced mistakes stay the same with more data.

\textbf{Which model the rest of this chapter uses.} Every other result in this chapter (the confusion matrix, the per-category results, the whole-system test of Section~\ref{sec:e2e}) and the model in Clausify come from the October 2026 model trained on 24{,}000 examples. The full-data run is a control experiment, run to test one possible explanation, and we report it as one instead of swapping in a slightly better number afterwards. Since the two cannot be told apart, nothing later would change if we had.
\subsection{A third source of variance: the split itself}
\label{sec:splitvar}
Sections~\ref{sec:budget} and~\ref{sec:seeds} tested two things that could move our numbers: the amount of training data and the random seed. Neither changed the conclusions. There is a third source that neither touches, and it may be the largest: \textbf{the split}. Every result in this project comes from one contract-level 80/20 split made with seed 42. With only 510 contracts, a different split puts different rare categories in the test set and could move the scores by itself. Seed variance does not capture this, because changing the training seed keeps the test set the same.

Retraining the transformer on twenty splits is too expensive, but the TF--IDF baseline retrains in under a minute. Since every comparison in this project is made against it, its spread across splits shows how much the comparison itself could move. We made \textbf{20 separate contract-level 80/20 splits} and retrained the baseline on each.

\begin{table}[!htbp]
\centering
\caption{Split variance}
\label{tab:splitvar}
\small
\begin{adjustbox}{max width=\textwidth}
\begin{tabular}{lccccc}
\toprule
 & \textbf{Mean} & \textbf{SD} & \textbf{Min} & \textbf{Max} & \textbf{Range}\\
\midrule
micro-F1 over 20 splits & 0.7829 & 0.0063 & 0.7702 & 0.7918 & 0.0216\\
\midrule
\multicolumn{6}{l}{\emph{The split used throughout this project (seed 42): 0.7747}}\\
\bottomrule
\end{tabular}
\end{adjustbox}
\end{table}

Three points follow. First, \textbf{the split is the largest of the three uncertainties we measured}: a standard deviation of 0.0063, against 0.0015 for the ensemble's seed variance and a bootstrap half-width of about 0.008 for the comparison. It is 1.5 times the ensemble's $+0.0042$ lead, so a different split could easily remove that lead. This is one more reason to treat the lead as unproven.

Second, \textbf{our split is not a lucky one}. At 0.7747 the baseline sits below the 20-split mean of 0.7829, in the lower part of the range, so if anything this project reports the baseline on a slightly unfavourable split. This matters because the opposite would have been a worry: a split that happened to favour the baseline would have made the transformer's loss look worse than it is.

Third, \textbf{the finding we rely on is far outside this variation}. The baseline's lead over the transformer is 0.0824, about thirteen split-level standard deviations. No realistic new split would reverse it.

What we cannot say anything about is the stability of \emph{individual categories} across splits, which is where it would matter most: a rare category with two test examples in one split may have eight in another, and its F1 means little in either case. Repeated splits or cross-validation for the whole system remain future work, and Section~\ref{sec:limitations} lists this as a limitation, not as a question we have answered.

\section{Comparisons and Relationships}
Table~\ref{tab:scorecard} brings together the test results of all three models. Every number was produced by the separate test notebook running on the saved models trained on the 80/20 split.

\begin{table}[!htbp]
\centering
\caption{Final test results}
\label{tab:scorecard}
\small
\begin{adjustbox}{max width=\textwidth}
\begin{tabular}{llll}
\toprule
\textbf{Stage} & \textbf{Model} & \textbf{Metric} & \textbf{Score}\\
\midrule
2A Presence & AND-ensemble (TF--IDF + DistilBERT) & micro-F1  & 0.779\\
2A Presence & AND-ensemble & Accuracy               & 85.65\%\\
2A Presence & AND-ensemble & Precision / Recall     & 77.38\% / 78.41\%\\
2A Presence & AND-ensemble & Specificity            & 89.10\%\\
2B Span     & DistilBERT-QA & token-F1 (given the window) & 0.764\\
2B Span     & DistilBERT-QA & token-F1 (window chosen by 2A) & 0.383\\
2B Span     & DistilBERT-QA & Exact Match           & 35.5\%\\
1 Summarizer & FLAN-T5-small & ROUGE-L (41 templates)  & 0.775\\
1 Summarizer & FLAN-T5-small & ROUGE-L (clause-specific) & 0.116\\
\midrule
\textbf{Full pipeline} & \textbf{all three steps} & \textbf{real clauses that get through} & \textbf{22.0\%}\\
\bottomrule
\end{tabular}
\end{adjustbox}
\end{table}

\textbf{How this compares with earlier work.} Table~\ref{tab:related} in Chapter~2 compares this project with earlier work. The original CUAD study reports span extraction with DeBERTa-xlarge, scored over whole documents with a precision-at-recall measure. Our small models are about ten times smaller, and our span scores are measured inside windows that contain the answer, so the numbers cannot be compared directly and we do not present them that way. What the comparison does show is a difference in scope: earlier work stops at finding clauses, while this project adds presence decisions, a risk level for each category and plain-English output, and tests the whole system from start to finish.

\textbf{How the parts relate.} Among the presence models, the TF--IDF baseline and the window-based transformer make different kinds of mistakes: the transformer finds more clauses but raises many false alarms (recall 0.913, precision 0.558), while the baseline is more careful (recall 0.810, precision 0.742). This is why the AND-ensemble, which needs both, matches the baseline while balancing precision and recall. Across the steps, the separate scores do not simply multiply: mistakes in the window chosen by the presence step carry over into span extraction, which is why the whole-system rate is far below the product of the separate scores.

\section{Discussions}
Putting the chapter together, the results fall into three groups, and it matters which is which.

\textbf{Findings that survived every check we could run.}
\begin{enumerate}[itemsep=2pt]
  \item For document-level \emph{presence}, a full-document lexical baseline beats a windowed DistilBERT by $+0.081$ micro-F1. We tested the two obvious alternative explanations and neither held: training on all 53{,}662 windows instead of 24{,}000 changed the score by $-0.0004$ (Section~\ref{sec:budget}), and across three seeds the gap is 9.5 standard deviations wide (Section~\ref{sec:seeds}). More data made the transformer fire \emph{more} rather than better, which is what max-pooling over 35 windows predicts. The limitation is the aggregation rule, not the encoder or the budget.
  \item For \emph{span extraction}, the transformer has no substitute, a bag-of-words model cannot point at character offsets, and it performs well in distilled form \emph{when handed the right window}: token-F1 0.763, 95\% CI $[0.740, 0.782]$.
  \item For \emph{summarization}, a 41-sentence template seed turns the task into classification. The fine-tuned model reproduces a template verbatim on 100\% of test clauses, and against clause-specific references it scores no better than the template alone (Section~\ref{sec:summ}). A high automatic metric should not be believed until the outputs have been read.
\end{enumerate}

\textbf{A claim we withdrew.} The AND-ensemble's first place in Table~\ref{tab:ensemble} is not a proven result. Its lead over the baseline is $+0.0042$. Seed noise alone (SD $0.0015$) would not wipe it out, but the bootstrap over contracts puts it inside $[-0.0037, +0.0115]$ (Section~\ref{sec:bootstrap}), and a different split moves the baseline by $0.0063$ (Section~\ref{sec:splitvar}). Two of the three checks say the top two rows of that table are a tie. What the ensemble does earn is balanced mistakes and about six times less seed variance than the transformer alone.

\textbf{Costs we did not see until we looked for them.} Two results in this chapter came from questions we had not thought to ask at first, and both are unflattering.
\begin{itemize}[itemsep=2pt]
  \item The AND rule is not good everywhere. On the ten rarest categories it gains \emph{nothing} in F1 over the transformer alone (0.364 against 0.368): it buys 34.6 points of precision with 30.0 points of recall and cuts the clauses found from 39 of 70 to 18 of 70 (Section~\ref{sec:raretail}). All of its overall benefit comes from the common categories, and one of the categories it scores zero on is High-risk.
  \item The separate scores do not multiply. Only 22.0\% of real clauses get through the whole system, against about 65\% predicted by multiplying the separate scores (Section~\ref{sec:e2e}), because the presence model's highest-scoring window is the wrong window 30.9\% of the time, and the span model was trained on centred windows it never sees in real use.
\end{itemize}

The common thread is about method, not about model design. Every number in this chapter that we first trusted least (the ensemble's rank, the summarizer's ROUGE, the span model's token-F1) turned out to measure something narrower than it seemed. In each case, the way to find out was to measure the thing the score stood for: resample the contracts, change the seed, change the reference sentences, run the steps one after another. None of these checks needed a bigger model or more data. They only needed the question to be asked.

\subsection{Usability}
\label{sec:usability}
No usability study was carried out. The app's main claim, that grouping clauses by risk helps a non-lawyer act on a contract, is therefore untested, and Section~\ref{sec:threats} lists it under construct validity instead of among the results.
%% ==================================================================
%% TODO IF THE INFORMAL WALKTHROUGH IS RUN (review/USABILITY_PROTOCOL.md)
%% Replace the paragraph above with 3-5 sentences, e.g.:
%%
%% "We ran an informal walkthrough with two non-lawyer participants (15
%%  minutes each). This is far too small to constitute a usability study and
%%  is reported only to replace an absence of evidence with a little.
%%  <What they did first.> <What they said in answer to the four questions.>
%%  <Anything suggesting over-trust in the quoted clause text or the
%%  confidence bar.> No task-completion or comprehension measurement was
%%  attempted, and with two participants no quantitative claim is possible."
%%
%% If a participant treated the quoted clause text as authoritative, or read
%% the confidence bar as a probability, say so -- Chapter~5 shows both are
%% unwarranted, and an observed instance belongs in the ethics discussion too.
%% ==================================================================
\subsection{Honest Limitations of the Demo}
\label{sec:demolimits}
By default the app finds clauses with the transformer presence model alone; the AND-ensemble is offered as a setting. As Section~\ref{sec:bootstrap} showed, the transformer alone reports too many clauses. Two measured facts explain why it is still the default.

\textbf{The confidence score is not a probability.} Section~\ref{sec:threshold} shows the transformer's score has an Expected Calibration Error of 0.205 and is too confident everywhere: among detections that score above 0.9, an almost full confidence bar, the clause is really there about 73\% of the time, and among those between 0.5 and 0.6, under 20\% of the time. So the app uses the fixed 0.5 threshold of every reported test (an earlier version had a slider for it, which invited reading the score as a probability) and labels every bar as a model score, not a probability. Fitting a calibration map on a validation split would fix the underlying problem without retraining.

\textbf{Using the transformer alone is, unusually, the better choice for High-risk clauses.} The project reports the AND-ensemble as its main setup, so a reader might think the app's use of the transformer alone is a compromise. Section~\ref{sec:riskeval} shows it is not, where it matters most: the ensemble misses 36.9\% of High-risk clauses, while the transformer alone misses 15.9\%. The app's setting is worse on overall micro-F1 and better at finding the clauses a signer would most regret missing. It started as a packaging accident, before we had run the risk-weighted test; it is now the chosen default. The ensemble is still available as an option, with one more warning measured on the app itself: on short contracts its TF--IDF half, which was trained on long SEC filings, blocks many detections. On the 9{,}174-character sample contract that ships with the app, the AND rule reports 6 clause types (1 High-risk), while the transformer alone reports 13 (2 High-risk).

\textbf{The 30-window limit cuts off long contracts without warning in the results.} To keep the app fast, it scores at most 30 windows, so any contract longer than 45{,}500 characters is only partly read, and clauses after that point cannot be found at all. The app says this in its notice, but none of the numbers in this chapter include it, because they score whole documents. This is not a rare case: \textbf{38.4\% of CUAD contracts (196 of 510) are longer than 45{,}500 characters}, and the longest test contract has 289{,}615 characters, of which the app reads the first 16\%. Table~\ref{tab:appbench} shows the effect: the 75th-percentile contract and the longest contract take the same time even though one is four times longer, because both are cut to 30 windows. The quoted text is also less reliable than the label. Section~\ref{sec:e2e} shows that when the presence model chooses the window, the span model's answer overlaps the real clause in only about a fifth of cases overall. The app now quotes the whole paragraph that the presence model scores highest, which fixed the clear failures we saw in our own tests, but this has not been measured on the test set. The practical message for a user is simple, and the app states it too: \textbf{the risk badge and the category name are much more reliable than the quoted text under them}. A card should be read as ``this contract seems to contain a clause of this type; look here'', not as ``this is the clause''. The risk levels are also set per category, not per clause: the risk table says a \emph{Cap on Liability} clause needs attention, but it does not check whether this particular cap is generous or harsh. All of this is stated in the notice the app shows above every analysis. What the demonstration does show is that the models work on contracts they have never seen, and that the output can be shown in a form a non-lawyer can act on.

